\documentclass[11pt]{article}
\usepackage[margin=1in]{geometry}
\usepackage{amsmath,amssymb,amsthm}
\usepackage{fontspec,xeCJK}
\setCJKmainfont{FandolSong-Regular.otf}
\usepackage[hidelinks]{hyperref}
\usepackage{enumitem}
\setlist[itemize]{nosep,leftmargin=*}
\newcommand{\F}{\mathbb F}
\newcommand{\R}{\mathbb R}
\newcommand{\N}{\mathbb N}
\newtheorem{theorem}{Theorem}
\title{Conjecture 00000002899:\\a counterexample to the displayed\\Kakeya asymptotic lower bound}
\author{C0ldSmi1e}
\date{12 October 2026}
\begin{document}
\maketitle

\noindent We give an explicit family of finite-field Kakeya sets and a Lean-checked
refutation of the displayed coefficient-one eventual lower bound in every
dimension at least two. The interpretation limit below is part of the result.

\section{Original statement and precise scope}
The original English statement is:
\begin{quote}
Definition: The Kakeya problem over finite fields: the size of sets containing
a line in every direction. Conjecture: The dimension spectrum of Kakeya sets:
the size of minimal Kakeya sets in d-dimensional space has the explicit
asymptotic lower bound of type q\textasciicircum d(1−c/q), attained by tube
merging of polynomial methods. (finite Kakeya size asymptotics)
\end{quote}
The original Chinese statement is:
\begin{quote}
定义：有限域上的 Kakeya 问题：含每方向线的集合的尺寸。猜想：Kakeya 集的维数谱：d 维空间的最小 Kakeya 集的尺寸有 q\textasciicircum d(1−c/q) 型显式渐近下界且下界由多项式方法的管合并达到。（有限 Kakeya 尺寸渐近）
\end{quote}
The source was captured from
\href{https://github.com/The-Last-Math-Competition/The-Last-Math-Competition/blob/1e3f7c25616e84b611350a412febceabad75b56a/conjectures/00000002899.md}{the repository's conjecture file}.

For a finite field $F$ of cardinality $q$ and an integer $d\geq 1$, a Kakeya set
$K\subseteq F^d$ contains a full affine line in every direction:
\[
 \forall v\in F^d\setminus\{0\},\quad
 \exists a\in F^d,\quad \forall t\in F,\quad a+tv\in K.
\]
The parametrization $t\mapsto a+tv$ is injective when $v\neq0$: a nonzero
coordinate of $v$ permits cancellation. Thus each required line has exactly
$q$ points. Rescaling a direction by a nonzero field element only
reparametrizes its line, so this definition includes every projective direction.
Write $m_d(F)$ for the minimum cardinality of a Kakeya subset of $F^d$.
This minimum exists and is attained because the full space is admissible and
the collection of subsets is finite.

We interpret the displayed expression as the eventual lower bound
\begin{equation}\label{eq:claim}
 \exists c\in\R\ \exists Q\in\N\ \forall F\text{ finite field},\qquad
 q\geq Q\ \Longrightarrow\ m_d(F)\geq q^d(1-c/q).
\end{equation}
Both $c$ and $Q$ may depend on the fixed dimension. We disprove
\eqref{eq:claim} for every $d\geq2$, allowing every real $c$ and every threshold.

\textbf{Interpretation limit.} The source leaves ``of type'', ``dimension
spectrum'', and ``tube merging'' undefined. Our result refutes the displayed
leading coefficient one with an error of order $q^{d-1}$. It does not refute
an unspecified positive multiple of $q^d$, such as an $\Omega(q^d)$ bound,
or a version allowing $c$ to grow with $q$. No unspecified multiplicative
constant is inserted into the displayed formula. Failure of its necessary
lower-bound clause suffices; we do not invent a definition for the additional
attainment method or claim to determine the sharp minimum.

\section{A family in every dimension at least two}
Let $F$ be a finite field with $2\neq0$ in $F$, and let $q=|F|$. Put
\[
 S=\{z^2:z\in F\},\qquad
 K_d=\{w\in F^d:w_0=0\text{ or }w_1+w_0^2\in S\},\quad d\geq2.
\]
Only the first two coordinates are constrained.

\begin{theorem}\label{thm:family}
The set $K_d$ is Kakeya, and
\begin{equation}\label{eq:card}
 |K_d|=\frac{q^2+2q-1}{2}\,q^{d-2}.
\end{equation}
Consequently
\begin{equation}\label{eq:density}
 \frac{m_d(F)}{q^d}\leq\frac12+\frac1q-\frac{1}{2q^2}.
\end{equation}
\end{theorem}
\begin{proof}
First count the squares, including zero. The fiber of $z\mapsto z^2$ over
zero consists of zero alone. For a nonzero square $z^2$, the fiber is
$\{z,-z\}$: equality of squares factors as
$(u-z)(u+z)=0$, and $z\neq-z$ because $2z\neq0$. Counting the fibers gives
\[
 q=1+2(|S|-1),\qquad 2|S|=q+1.
\]

Now fix any nonzero direction $v\in F^d$. If $v_0=0$, choose $a=0$; every
point $a+tv$ has first coordinate zero and belongs to $K_d$. This case also
covers directions whose first two coordinates vanish but a later one does not.
If $v_0\neq0$, set $b=v_1/(2v_0)$, and let $a_1=b^2$ with every other
coordinate of $a$ equal to zero. For every $t\in F$,
\[
 (a+tv)_1+(a+tv)_0^2
   =b^2+tv_1+(tv_0)^2=(b+tv_0)^2,
\]
since $2bv_0=v_1$. Thus the whole line lies in $K_d$ in this case as well.

For a fixed first coordinate $x=0$, all $q$ second coordinates are allowed.
For each $x\neq0$, translation $y\mapsto y+x^2$ is a bijection, so exactly
$|S|$ second coordinates are allowed. The remaining $d-2$ coordinates are
free. Hence
\[
 |K_d|=\bigl(q+(q-1)|S|\bigr)q^{d-2}
       =\frac{q^2+2q-1}{2}\,q^{d-2}.
\]
This is a count of the explicitly defined family, not an optimality claim.
The genuine minimum satisfies $m_d(F)\leq|K_d|$. Dividing by the positive
number $q^d$ proves \eqref{eq:density}.
\end{proof}

\section{Counterexamples beyond every threshold}
\begin{theorem}\label{thm:unbounded}
For every integer $d\geq2$, every $c\in\R$, and every $Q\in\N$, there is an
odd prime $p\geq Q$ and a Kakeya set $K\subseteq(\mathbb Z/p\mathbb Z)^d$ such that
\[
 m_d(\mathbb Z/p\mathbb Z)\leq |K|<p^d(1-c/p).
\]
In particular, \eqref{eq:claim} is false in every such dimension.
\end{theorem}
\begin{proof}
Choose a natural number $N\geq Q$ with $N>2$ and $N>2(c+1)$, using the
Archimedean property of $\R$. Unboundedness of primes gives a prime $p\geq N$.
Thus $p>2$, so $F=\mathbb Z/p\mathbb Z$ is an actual odd finite field of
cardinality $p$, and the family of Theorem~\ref{thm:family} is available.
For this $p$,
\[
 p^2(1-c/p)-\frac{p^2+2p-1}{2}
   =p\left(\frac p2-c-1\right)+\frac12>0.
\]
Multiplication by $p^{d-2}>0$, followed by the minimum comparison, gives the
stated strict inequality. The choices were made after arbitrary $c$ and $Q$;
this proves an unbounded family of violations, not a small-field exception.
If \eqref{eq:claim} held, specialize its field quantifier to this prime field
to contradict the strict inequality.
\end{proof}

\section{Formal correspondence and verification}
The Lean project uses \texttt{Fin d $\to$ F} as its actual coordinate space.
\texttt{IsKakeya} requires membership of every point on a full affine line
for every nonzero vector. \texttt{family} constrains coordinates zero and one
and leaves all others free. Its explicit subtype equivalence to a first-two-
coordinate fiber times \texttt{Fin n $\to$ F} supplies the higher-dimensional
cardinality argument.

\texttt{minimumKakeyaSize} is \texttt{Nat.find} over cardinalities of all
admissible Kakeya finsets, with existence witnessed by the full space.
Attainment and comparison with any actual Kakeya finset are separate theorems.
The natural counting identity used before casting to $\R$ is
\[
 2|K_{n+2}|+q^n=(q^2+2q)q^n,
\]
which avoids truncated subtraction and division in natural numbers.

\texttt{ProposedEventualLowerBound} quantifies over arbitrary real $c$,
natural $Q$, and every finite field. Prime fields suffice to contradict this
universal assertion; their unbounded sizes are proved using the generic
unbounded-prime theorem. All dimensions $d\geq2$ are obtained by $d=(d-2)+2$.

\subsection*{Checked project and conclusions}
The eight source modules are the default root \texttt{Conjecture2899},
\texttt{Geometry}, \texttt{Squares}, \texttt{Counting}, \texttt{Minimum},
\texttt{Bounds}, \texttt{Asymptotic}, and \texttt{AxiomAudit}; the latter seven
are under the \texttt{Conjecture2899} namespace directory. The default target
reaches them all. The principal declarations in that namespace are
\begin{itemize}
\item \texttt{counterexamples\_in\_each\_dimension}: Theorem~\ref{thm:unbounded},
with the actual field, set, minimum comparison, and strict inequality;
\item \texttt{not\_proposedEventualLowerBound}: the explicit negation of
\eqref{eq:claim} for every $d\geq2$;
\item \texttt{normalized\_minimum\_upper\_bound\_in\_each\_dimension}:
\eqref{eq:density} for every finite field with $2\neq0$ and every $d\geq2$.
\end{itemize}
The global declaration \texttt{finite\_field\_square\_count} proves the square
fiber count. The complete inventory lists 25 named definitions and theorems.

The project is pinned to Lean 4.19.0 and Mathlib v4.19.0, at commit
\begin{center}
\small\texttt{c44e0c8ee63ca166450922a373c7409c5d26b00b}.
\end{center}
All nine dependency revisions are locked. An independent source-only project copy passed the
complete Lake build and strict replay of every source module and the Lake
configuration with warnings treated as errors. The 11 frozen source and
configuration files remained identical. All 39 owned declarations, including
14 generated proof or equation theorems, were audited by originating module.
Their axiom dependencies are limited to \texttt{propext},
\texttt{Classical.choice}, and \texttt{Quot.sound}; no unsafe generated helper
is present. The audit covers 8,326 reachable declarations and verifies 2,942
imported module identities. All 63 execution receipts were checked independently.

An ordinary safe control was accepted, while five controls containing a custom
axiom, a transitively hidden axiom, an admitted proof, compiler-trusted decision,
or unsafe definition were rejected as intended. There is no admitted proof,
custom axiom, or compiler-trusted decision in the submission. The square count
and prime-field counterexamples are symbolic proofs; no external numerical
program is used as a proof premise. Separate semantic review checked the actual
objects, quantifiers, minimum, and interpretation against both original texts.

\textbf{Trust and acceptance limits.} Generic dependency source identities and
cached imported artifacts were checked, but the entire dependency library was
not rebuilt and Lean was not bootstrapped. The dependency-graph check is not a
second proof kernel; source compilation supplies Lean kernel checking.
Runtime binaries and libraries were hashed; imported runtime artifacts were
recorded without a compiler bootstrap. PATH safeguards are not operating-system
network isolation. The accompanying verification record supplies reproduction
commands, raw receipts, independent reviews, and the final PDF inspection.
These local checks do not constitute maintainer acceptance or settle alternative
interpretations of the informal source.
\end{document}
